\documentclass{article}
\usepackage{graphicx} % Required for inserting images
\usepackage[a4paper,margin=2.5cm]{geometry}

\title{WBA_Aufgabe2}
\author{michael-kolle }
\date{May 2026}

\begin{document}



\section{User Storys}

\begin{table}[h!]
\centering
\begin{tabular}{|c|p{2.5cm}|p{4cm}|p{4cm}|p{2cm}|}
 \hline
 Story & Wer & Möchte Was & Warum & Titel \\ [0.5ex] 
 \hline\hline
1 & Projektleiter & Die Gesamtarbeitszeit meiner Projekte sehen & Damit ich die Zeitplanung verbessern kann & Titel? \\ 
2 & Projektleiter & Die Aufgaben nach verbleibender Arbeitszeit sortieren & Die Abschätzung für den Aufwand zu vereinfachen & Titel? \\
3 & Projektleiter & Am ende einer erledigten Aufgabe die tatsächlich benötige Zeit einsehen können & Falls nötig den Zeitplan anpassen und für zukünftige Projekte zu lernen & ??? \\
 4 & Teammitglied & Die einzelnen Aufgaben in einem Projekt nach verbleibender Zeit sortieren & Eine bessere Planung im Arbeitstag zu erzielen & ? \\
 5 & Teammitglied & die Aufgaben vor und nach einer Aufgaben sehen, wenn diese abhängig sind & Damit eine Aufgabe nicht zum bottleneck wird & ? \\
 6 & Projektleiter & Eine Übersicht über alle Projekt haben & Um schnell zwischen Projekten zu wechseln & ? \\
 7 & Projektleiter & Bei der Planung eines Projekts aus verschiedenen Arbeitsschritten wählen & Um die Planung zu erleichtern & ? \\ 
 8 & Projektleiter & Meine Projekte Verwalten können (Löschen, Bearbeiten) & Um im späteren Verlauf flexibel zu sein & ? \\ [1ex] 
 \hline
 \end{tabular}
\end{table}

\section{Mockups}
Ipad

\section{Aufwandsabschätzung}

\subsection{Aufagbe a}
Das Prjekt wird in kleinere Aufgaben geteilt
Die Aufgaben werden nach Schwierigkeit, Aufwand, Art der Aufgabe (Design, Programmieren) gewichtet und eine geschätzte Zahl an Arbeitsstunden festgelegt.

Die Gesamtzeit ist dann die Zeit der Aufgaben die wir schätzen plus einen kleinen Puffer

\subsection{Aufgabe b}
Als Zahlen nehmen wir: 
Anzahl der Seiten für die User Story
Anzahl der Funktionen
Geschätzte Zeit pro Funktion
Zeit für das Design
Puffer Zeit für Fehlerbehebung

\subsection{Aufgabe c}
Startseite: 4 Stunden
Projektübersicht: 5 Stunden
Projektdetailseite: 7 Stunden
Seite für ein neues Projekt: 7 Stunden
Filter und Soriterung: 4 Stunden


\end{document}